
Applying k-means clustering (k=2) to epoch-5 penultimate-layer embeddings of a ResNet-50 pretrained on ImageNet, trained for 5 epochs on Waterbirds, and separately trained for 5 epochs on CelebA, produces an AMI of 0.762 against the Waterbirds ground-truth group labels and an AMI of 0.258 against CelebA ground-truth group labels. These two experiments use the same model architecture, the same optimizer, the same number of training epochs, the same clustering algorithm, the same hyperparameters, and the same random seed. No per-dataset tuning was applied. The threefold gap in AMI (0.762 versus 0.258) cannot be attributed to experimental variation.

This observation is consequential for annotation-free spurious correlation robustification. Methods such as GEORGE \citep{Sohonietal2020} and PruSC \citep{Kimetal2024} cluster penultimate-layer ERM embeddings to discover spurious groups, then use the cluster assignments as pseudo-group-labels for downstream reweighting or neuron pruning. If clustering fails — returning near-random group assignments — any subsequent robustification is applied to incorrect directions. Because the practitioner chose the annotation-free path specifically to avoid ground-truth group labels, there is no readily available signal indicating that the clustering step has failed.

The broader problem is that existing annotation-free detection methods report clustering results on benchmarks where detection succeeds, without characterizing the conditions under which it fails. GroupDRO \citep{Sagawaetal2020}, JTT \citep{Liuetal2021}, and DFR \citep{Izmailovetal2022} are widely used robustification approaches; annotation-free variants inherit an untested assumption about the reliability of the clustering step. Published annotation-free clustering results appear to be obtained at or near training convergence, with k chosen to match or exceed the number of known groups. The early-epoch (epoch-5), k=2 configuration used by intervention methods is a distinct experimental setting, and published convergence-epoch results with flexible k do not establish that early-epoch, k=2 detection is reliable. For CelebA under the epoch-5, k=2 setting, the obtained purity of 0.456 is barely above the random baseline of 0.440.

The gap between Waterbirds and CelebA results is not a replication failure. It is a consequence of a previously uncharacterized constraint on clustering-based annotation-free detection: the success of epoch-5 k-means clustering depends on whether the spurious attribute is already encoded as a salient, separable dimension in the pretrained model's representation. ImageNet-pretrained ResNet-50 encodes scene-level features — bird habitats, backgrounds, and environmental contexts — that align directly with the Waterbirds spurious correlation (bird type versus background). Hair color is not a semantically distinct primary category in ImageNet; the pretrained weights do not pre-separate CelebA spurious groups along a hair-color axis, and five epochs of ERM fine-tuning are insufficient to construct this structure.

This paper investigates the following questions: (1) whether epoch-5 k-means clustering reliably recovers spurious group structure when the spurious attribute is scene-level and aligns with the ImageNet pretraining prior; (2) whether the same procedure fails when the spurious attribute is a fine-grained texture not well-represented in the pretraining prior; (3) whether the observed gap reflects a genuine difference in detection reliability; and (4) what practical implications follow for the deployment of annotation-free spurious detection pipelines.

The contributions of this work are:

\textbf{C1 — Characterization of feature-strength conditionality:} A systematic characterization of when annotation-free clustering-based spurious direction discovery produces above-chance cluster-group alignment versus near-random performance, with pretrained initialization alignment identified as the primary differentiating factor.

\textbf{C2 — Empirical quantification of the conditionality gap:} Measurement of AMI=0.762 and purity=0.892 for Waterbirds against AMI=0.258 and purity=0.456 for CelebA under identical experimental conditions, along with random baseline measurements (Waterbirds random AMI≈0.0001, purity=0.728; CelebA random AMI≈0.000, purity=0.440) that establish the reference floor.

\textbf{C3 — Proposed diagnostic:} A pretrained linear probe pre-screen is proposed as a candidate diagnostic to assess whether a spurious attribute is encoded in the pretrained representation prior to fine-tuning. This diagnostic is motivated but is not empirically calibrated in this work; the appropriate threshold requires future validation.

\textbf{C4 — Identification of configuration-sensitivity reporting gap:} The analysis demonstrates that annotation-free clustering results are highly sensitive to probe epoch and choice of k, and that convergence-epoch results with k>2 do not establish the reliability of early-epoch, k=2 detection; this configuration difference is not systematically disclosed in existing reporting.

The downstream Gradient SNR Balancing (GSB) mechanism — which proposed equalizing gradient SNR along cluster-discovered directions to improve worst-group accuracy — was not evaluated because the prerequisite detection condition (H-E1) failed on CelebA, blocking all downstream experiments (H-M1 through H-C1). Results reported in this paper concern the detection step only.

The remainder of the paper is organized as follows. Section 2 reviews related work on spurious correlation robustification and annotation-free detection methods. Section 3 describes the experimental methodology. Section 4 presents the experimental setup. Section 5 reports results. Section 6 discusses implications and limitations. Section 7 concludes.

